\documentclass[11pt,a4paper]{article}
\usepackage[utf8]{inputenc}
\usepackage[ngerman]{babel}
\usepackage{amsmath,amssymb,amsthm}
\usepackage{geometry}
\geometry{margin=1.15cm}
\usepackage{setspace}
\setstretch{1.25}
\usepackage{booktabs}
\usepackage{enumitem}
\usepackage{hyperref}
\hypersetup{colorlinks=true,linkcolor=blue,citecolor=blue,urlcolor=blue}

\newtheorem{theorem}{Satz}
\newtheorem{proposition}{Proposition}
\newtheorem{lemma}{Lemma}
\newtheorem{corollary}{Korollar}
\theoremstyle{definition}
\newtheorem{postulate}{Postulat}
\newtheorem{definition}{Definition}
\theoremstyle{remark}
\newtheorem{remark}{Bemerkung}

\newcommand{\ag}{a_\gamma}
\newcommand{\CQ}{C_Q}
\newcommand{\Ok}{\Omega_K}
\newcommand{\OL}{\Omega_\Lambda}
\newcommand{\Om}{\Omega_m}
\newcommand{\Mpl}{M_{\mathrm{Pl}}}
\newcommand{\avg}[1]{\langle #1\rangle}

\title{Eine parameterfreie Beziehung zwischen räumlicher Krümmung\\
und der kosmologischen Konstante,\\
hergeleitet aus der konformen Anomalie des Photons}
\author{Lukas Röhl}
\date{\today}

\begin{document}
\maketitle

\begin{abstract}
Dies ist eine in sich geschlossene Darstellung für eine mathematisch geschulte Leserin, von der
\emph{nicht} vorausgesetzt wird, dass sie Quantenfeldtheorie in gekrümmter Raumzeit kennt. Jeder
Begriff wird definiert und jeder Schritt soweit wie möglich aus ersten Prinzipien hergeleitet; die
eine Stelle, an der eine Herleitung auf eine echte physikalische Hypothese trifft, wird isoliert und
als solche gekennzeichnet. Das zentrale Ergebnis ist eine einzige dimensionslose Zahl, die die
räumliche Krümmungsdichte $\Ok$ und die Dunkle-Energie-Dichte $\OL$ des Universums verbindet,
\[
\frac{\lvert\Ok\rvert}{\OL}=\frac{\ag}{8\pi^2}=\frac{31}{1440\pi^2}\approx 2.181\times10^{-3},
\]
wobei $\ag=31/180$ der Typ-A-Koeffizient der konformen Anomalie des freien elektromagnetischen
Feldes ist. Wir leiten $\ag$ aus der Gilkey-Wärmekern-Rechnung her (\S\ref{sec:heat}); den
topologischen Vorfaktor aus Gauß--Bonnet--Chern (\S\ref{sec:gbc}); die Krümmungsladung
$\CQ=8\pi^2$ aus einem elementaren Integral (\S\ref{sec:Q}); den spektralen Unterbau $R1$--$R12$
mit zwei unabhängigen Auswertungen des $\zeta$-Wertes auf der geschlossenen Mannigfaltigkeit
(\S\ref{sec:spectral}); die Unterdrückung der naiven Vakuumenergie durch Sequestrierung
(\S\ref{sec:seq}); den Lift auf die kosmologische Skala durch kohärente spektrale Akkumulation
(\S\ref{sec:lift}); die integrierte Invariante $4\ag$ über zwei unabhängige Wege
(\S\ref{sec:invariant}); das Vorzeichen der Krümmung aus einem Variationsprinzip auf der
Hemisphäre (\S\ref{sec:cap}); und die Eindeutigkeit des Photons unter den Kandidatenfeldern
(\S\ref{sec:uniq}). Die Friedmann-Übersetzung (\S\ref{sec:friedmann}) liefert $\OL=4\ag=0.6889$,
$\Ok\simeq+1.5\times10^{-3}$, $w=-1$. Die einzige physikalische Hypothese steht in
\S\ref{sec:postulate}; \S\ref{sec:ledger} listet den logischen Status jeder Aussage auf. Ein
begleitendes Skript (\texttt{csg\_kp\_falsifikation.py}) berechnet jede Zahl aus ihren Eingaben.
\end{abstract}

\section*{Prolog --- Die Krümmung, die das Licht verrät}
\textit{Eine kurze, laienverständliche Einleitung.}

\medskip
Niemand hat den Raum je gesehen. Wir sehen die Dinge \emph{im} Raum --- Sterne, Galaxien, das
Licht, das von ihnen zu uns kommt. Aber die Form des Raumes selbst können wir nur erschließen. Ist
er flach wie ein Tisch? Oder krümmt er sich, so wie sich die Oberfläche der Erde krümmt, nur dass
wir mitten darin sitzen und nichts davon merken?

Diese harmlose Frage führt schnell in eine Verlegenheit. Unsere beste Theorie sagt, das leere
Vakuum müsse eine Energie tragen --- und zwar eine, die um einen Faktor mit hundertzweiundzwanzig
Nullen zu groß ist. Es ist die schlechteste Vorhersage, die die Physik je gemacht hat. Seit über
hundert Jahren hat niemand sie reparieren können.

Diese Arbeit repariert sie nicht. Sie folgt stattdessen einer feinen Spur, und die Spur liegt im
Licht.

Im flachen, leeren Raum ist das Licht gleichgültig gegen jeden Maßstab; es kennt kein Groß und kein
Klein. Krümmt man den Raum aber, geht diese Gleichgültigkeit verloren --- die Quantentheorie
hinterlässt einen winzigen Rest. Das Schöne daran: Dieser Rest ist keine Sache der Wahl. Er ist
eine feste Zahl, so unausweichlich wie das $\pi$ im Umfang eines Kreises. Für das Licht lautet sie
einunddreißig Hundertachtzigstel.

Und nun die Wette dieser Arbeit: Dieselbe kleine Zahl, die das Licht im gekrümmten Raum verrät,
bestimmt zugleich, wie stark sich das ganze Universum krümmt. Rechnet man es aus, kommt eine
Vorhersage heraus --- etwa zwei Tausendstel, und der Raum krümmt sich nicht nach innen, sondern ein
winziges Stück nach außen. Kein Regler, an dem man drehen könnte. Eine Zahl, fertig.

Es gibt einen Haken, und die Arbeit verschweigt ihn nicht. Die Zahl selbst --- das nackte
Verhältnis zwischen der Krümmung des Universums und der Dunklen Energie, die es füllt --- folgt als
Satz; sie verdankt nichts einer Wahl. Was an einer einzigen Annahme hängt, ist die Brücke: dass
diese zeitlose Zahl sich treu in das Universum überträgt, in dem wir wirklich leben, und dort die
absolute Menge der Dunklen Energie festlegt. Diese Annahme ist nicht bewiesen, sie ist umstritten,
sie wird offen hingestellt, und die empirische Aussage ruht sichtbar auf ihr. Das ist die ganze
Gestalt der Arbeit: eine Zahl (ein Satz), eine Annahme (die Brücke), ein Test.

Denn das Beste an einer Vorhersage ist, dass man an ihr scheitern kann. In den nächsten Jahren
werden zwei große Himmelsdurchmusterungen, DESI und Euclid, die Krümmung des Raumes genauer
vermessen als je zuvor --- genau jene Zahl. Stimmt sie nicht, ist die Idee tot. Nicht nachgebessert,
nicht umgedeutet: tot. Diese Bereitschaft zu sterben ist kein Mangel. Sie ist das Einzige, was eine
Idee ehrlich macht.

Was hier beginnt, ist also keine fertige Antwort, sondern ein sauber gestellter Einsatz. Trägt der
Himmel die Zahl, dann hat sich eine erstaunlich schlichte Verbindung als wahr erwiesen --- zwischen
dem Kleinsten, einer Macke im Licht, und dem Größten, der Form des Alls. Und wenn nicht, dann wissen
wir es bald. Das ist fast ebenso viel wert.


\tableofcontents

\section{Was eine Lösung leisten muss}\label{sec:problem}
Das Problem der kosmologischen Konstante hat zwei Gesichter. Das \emph{Hierarchie}-Gesicht fragt,
warum die beobachtete Vakuumenergie nicht von der Größenordnung $\Mpl^4$ ist --- eine naive
Erwartung, die um einen Faktor $\sim10^{122}$ daneben liegt. Das \emph{Wert}-Gesicht fragt, warum
der Dunkle-Energie-Anteil gerade die Zahl $\OL\approx0.69$ ist, weder $0$ noch $1$. Eine
vollständige Lösung sollte (i) die Unterdrückung ohne Feinabstimmung erklären, (ii) den Wert
parameterfrei erzeugen und (iii) eine falsifizierbare Vorhersage machen, die sie von einer reinen
Konstante unterscheidet.

Wir versuchen \emph{keine} Quantentheorie der Gravitation. Wir stellen eine engere, schärfere
Frage: \emph{Gibt es eine reine Zahl, berechenbar allein aus bekanntem Feldinhalt, der das
Verhältnis von Krümmung zu Dunkler Energie des Universums gleichen muss?} Der Hauptteil dieser
Arbeit bejaht das und leitet diese Zahl her. Eine physikalische Hypothese ist nötig, um den
bewiesenen euklidischen Quotienten in den heutigen Absolutwert zu überführen; sie ist in
\S\ref{sec:postulate} ausdrücklich gekennzeichnet und der eigentliche Angriffspunkt jedes
Falsifizierungsversuchs.

\paragraph{Fahrplan in einfacher Sprache.} Ein Feld, das klassisch blind gegen den Gesamtmaßstab
der Abstände ist (``konform invariant''), \emph{verliert} diese Blindheit generisch bei der
Quantisierung. Der Rest ist eine geometrische, zahlenwertige Größe, die \emph{konforme
(Spur-)Anomalie}. Für das Photon wird sie von einem rationalen Koeffizienten $\ag$ kontrolliert.
Auf der runden Vierersphäre $S^4$ --- der euklidischen Form des de-Sitter-Raums --- integriert
sich diese Anomalie zu einer topologischen Invariante. Wir berechnen diese Invariante, teilen
durch eine zweite, elementare geometrische Invariante der Hemisphäre und erhalten eine reine Zahl.
Diese Zahl ist das behauptete Verhältnis von Krümmung zu Dunkler Energie.

\section{Die Spuranomalie aus ersten Prinzipien}\label{sec:anomaly}
\begin{definition}[Weyl-Reskalierung, konforme Invarianz]
Eine \emph{Weyl-Reskalierung} einer Metrik ist $g_{\mu\nu}(x)\mapsto e^{2\sigma(x)}g_{\mu\nu}(x)$
für eine glatte Funktion $\sigma$. Eine Wirkung $S[g,\phi]$ ist \emph{konform invariant}, wenn sie
unter einer Weyl-Reskalierung samt passender Reskalierung der Felder $\phi$ unverändert bleibt. Die
freie Maxwell-Wirkung $S=-\tfrac14\int\sqrt g\,F_{\mu\nu}F^{\mu\nu}$ in vier Dimensionen ist konform
invariant; folglich ist ihr klassischer Energie-Impuls-Tensor spurfrei: $T^\mu{}_\mu=0$.
\end{definition}

\begin{definition}[Spuranomalie und Typ-A-Koeffizient]
Das Pfadintegralmaß des quantisierten Feldes ist nicht Weyl-invariant; die regularisierte
Ein-Schleifen-Wirkung $\Gamma[g]$ erhält eine Abhängigkeit vom Maßstab $\sigma$. Die Spur des
quantenmechanischen Energie-Impuls-Tensors ist dann von null verschieden,
\begin{equation}\label{eq:anomaly}
\avg{T^\mu{}_\mu}=\frac{1}{16\pi^2}\bigl(c\,W^2-\ag\,E_4\bigr)+\nabla_\mu J^\mu,
\end{equation}
mit dem Weyl-Tensor $W_{\mu\nu\rho\sigma}$, der Euler-(Pfaffian-)Dichte
$E_4=R_{\mu\nu\rho\sigma}R^{\mu\nu\rho\sigma}-4R_{\mu\nu}R^{\mu\nu}+R^2$ und einem
schemaabhängigen Strom $J^\mu$. Der Koeffizient $c$ multipliziert die konform invariante Größe
$W^2$ (``Typ B''), $\ag$ die topologische Größe $E_4$ (``Typ A''). Beides sind reine Zahlen, die
durch den Feldinhalt festgelegt sind \cite{DeserSchwimmer,Duff1994}.
\end{definition}

\begin{remark}\label{rem:why}
Zwei strukturelle Tatsachen tragen alles Folgende. \emph{(a)} Auf einer konform flachen
Mannigfaltigkeit (wie $S^4$) verschwindet der Weyl-Tensor, $W=0$, sodass nur der Typ-A-Term
beiträgt. \emph{(b)} Das Integral $\int_M E_4$ über eine geschlossene Mannigfaltigkeit ist eine
topologische Invariante (Satz~\ref{thm:gbc}). Auf $S^4$ ist die integrierte Anomalie daher eine
reine Zahl mal $\ag$, frei von jeder lokalen Schema-Mehrdeutigkeit.
\end{remark}

\paragraph{Warum nur das Photon.} Ein Feld der Masse $m$ trägt zur Niederenergie-Vakuumantwort über
Terme bei, die mit Potenzen von $m^2/\Mpl^2$ unterdrückt sind, und entkoppelt im Infraroten. Das
Photon ist exakt masselos --- es ist das ungebrochene $U(1)$-Eichboson des Elektromagnetismus ---,
also exakt konform gekoppelt, und seine Typ-A-Anomalie ist der saubere, unverdrängte Beitrag in der
benötigten Ordnung. Es ist das einzige eingehende Feld; die Eindeutigkeit dieser Wahl wird in
\S\ref{sec:uniq} begründet.

\section{Der Koeffizient $\ag=31/180$ aus dem Wärmekern}\label{sec:heat}
\begin{theorem}[Typ-A-Koeffizient des Photons]\label{thm:agamma}
Der Typ-A-Koeffizient der konformen Anomalie des freien Maxwell-Feldes ist $\ag=\dfrac{31}{180}$.
\end{theorem}
\begin{proof}
Für einen elliptischen Operator zweiter Ordnung $D=-\Box+E$ auf Schnitten eines Vektorbündels über
einer Viermannigfaltigkeit hat der Wärmekern die DeWitt--Schwinger-Kurzzeitentwicklung
$K(x,x;s)=(4\pi s)^{-2}\bigl[\mathbf 1+s\,a_1+s^2 a_2+\cdots\bigr]$, und der integrierte vierte
Seeley--DeWitt-Koeffizient ist durch Gilkeys Formel gegeben \cite{Gilkey}
\begin{equation}\label{eq:gilkey}
\int a_4=\frac{1}{(4\pi)^2\,360}\int\!\sqrt g\;\mathrm{tr}\Bigl\{60RE+180E^2+30\,\Omega_{\mu\nu}\Omega^{\mu\nu}
+\mathbf 1\bigl(5R^2-2R_{\mu\nu}^2+2R_{\mu\nu\rho\sigma}^2\bigr)\Bigr\},
\end{equation}
wobei $\Omega_{\mu\nu}$ die Krümmung des Zusammenhangs auf dem Bündel ist. Die konforme Anomalie
eines einzelnen Feldes ist gleich $\zeta(0)$, also dem integrierten $a_4$; auf der Einheits-$S^4$
gelten die Einstein-Raum-Identitäten $R_{\mu\nu}^2=R^2/4$, $R_{\mu\nu\rho\sigma}^2=R^2/6$, $W^2=0$
und $\int E_4=64\pi^2$ (Satz~\ref{thm:gbc}), sodass nur der Euler-Anteil überlebt und ein einzelnes
Feld $\zeta(0)=-4a$ beiträgt, wenn $a$ sein Typ-A-Koeffizient ist. Wir werten drei Operatoren aus.

\smallskip
\emph{(i) Konform gekoppeltes Skalarfeld}, $D=-\Box+\tfrac R6$ ($E=-\tfrac R6$, $\dim 1$,
$\Omega=0$). Die Klammer in \eqref{eq:gilkey} reduziert sich auf $-\tfrac16 R^2$ und ergibt
$\zeta(0)_{\mathrm{Skalar}}=-\tfrac1{90}$, also $a_{\mathrm{Skalar}}=\tfrac1{360}$ --- der
Lehrbuchwert. Er dient als Normierungskontrolle der gesamten Rechnung.

\emph{(ii) Hodge--de-Rham-Laplace-Operator auf $1$-Formen}, mit Weitzenböck-Potential
$E=-\mathrm{Ric}$ (also $\mathrm{tr}\,E=-R$, $\mathrm{tr}\,E^2=R^2/4$, $\dim 4$,
$\mathrm{tr}\,\Omega_{\mu\nu}\Omega^{\mu\nu}=-R_{\mu\nu\rho\sigma}^2=-R^2/6$). Einsetzen ergibt
\[
\zeta(0)_{\Delta_1}=\int\avg{T}_{\Delta_1}=-\tfrac{2}{45}.
\]

\emph{(iii) Minimales Skalar-Geistfeld}, $E=0$, $\dim 1$:
$\zeta(0)_{\Delta_0}=\int\avg{T}_{\Delta_0}=\tfrac{29}{90}$.

\smallskip
Die Photon-Zustandssumme wird eichfixiert mit zwei Faddeev--Popov-Skalar-Geistern, die die Eich-
und Längsmoden entfernen, $\ln Z=-\tfrac12\ln\det\Delta_1+\ln\det\Delta_0^{\mathrm{Geist}}$, sodass
ihre Anomalie die Operatorkombination $\Delta_1-2\Delta_0$ ist:
\begin{equation}\label{eq:zetacombo}
\zeta(0;\mathrm{Maxwell};S^4)=\zeta(0)_{\Delta_1}-2\,\zeta(0)_{\Delta_0}
=-\tfrac{2}{45}-2\cdot\tfrac{29}{90}=-\tfrac{4}{90}-\tfrac{58}{90}=-\tfrac{62}{90}=-\tfrac{31}{45}.
\end{equation}
Da ein einzelnes Feld $\zeta(0)=-4a$ liefert, lesen wir
$\ag=-\tfrac14\bigl(-\tfrac{31}{45}\bigr)=\dfrac{31}{180}$ ab.
\end{proof}

\begin{remark}
$\ag$ ist somit \emph{berechnet}, nicht aus einer Quelle übernommen. Der nebenbei gewonnene
Skalarwert $1/360$ ist eine unabhängige Kontrolle der Normierung von \eqref{eq:gilkey}.
\end{remark}

\section{Der topologische Vorfaktor: Gauß--Bonnet--Chern}\label{sec:gbc}
\begin{theorem}[Gauß--Bonnet--Chern; der Vorfaktor $4=2\chi(S^4)$]\label{thm:gbc}
Für eine geschlossene orientierte Riemannsche Viermannigfaltigkeit $M$ gilt
$\displaystyle\int_M E_4\,\sqrt g\,d^4x=32\pi^2\,\chi(M)$, mit der Euler-Charakteristik $\chi$. Auf
$S^4$ ist $\chi=2$, also $\int_{S^4}E_4=64\pi^2$, und mit der konformen Normierung $16\pi^2$ in
\eqref{eq:anomaly},
\[
\int_{S^4}\avg{T^\mu{}_\mu}=\frac{\ag}{16\pi^2}\int_{S^4}E_4=\frac{\ag}{16\pi^2}\cdot 64\pi^2=4\ag,
\qquad
\frac{\int_{S^4}E_4}{16\pi^2}=4=2\,\chi(S^4).
\]
Der Vorfaktor $4$ ist der Atiyah--Singer-Index für die Euler-Charakteristik; er ist geometrisch und
kein freier Parameter. \cite{Chern,AtiyahSinger}
\end{theorem}

\section{Die Krümmungsladung der Hemisphäre $\CQ=8\pi^2$}\label{sec:Q}
\begin{definition}[$Q$-Krümmung, Paneitz-Operator]
Auf einer Viermannigfaltigkeit ist die Branson-$Q$-Krümmung der lokale Skalar
$Q_4=-\tfrac16\Box R-\tfrac12 R_{\mu\nu}R^{\mu\nu}+\tfrac16 R^2$ (bis auf Konvention), das
natürliche Analogon vierter Ordnung zur Gaußschen Krümmung; sie transformiert unter Weyl-Reskalierung
durch den Paneitz-Operator $P_4$ vierter Ordnung. Auf einer Einstein-Mannigfaltigkeit ist $\Box R=0$
und $Q_4$ reduziert sich auf $\tfrac16(R^2-3R_{\mu\nu}R^{\mu\nu})$.
\end{definition}

\begin{lemma}[$Q$-Ladung der Hemisphäre]\label{lem:qcharge}
Auf der runden Einheits-Hemisphäre $D^4=\{0\le\theta\le\pi/2\}\subset S^4$ gilt
$\displaystyle\CQ\equiv\int_{D^4}Q_4\,\sqrt g\,d^4x=8\pi^2$.
\end{lemma}
\begin{proof}
Auf der Einheits-$S^4$ (Einstein) ist $R=12$, $R_{\mu\nu}=3g_{\mu\nu}$, also
$R_{\mu\nu}R^{\mu\nu}=36$, und $\Box R=0$; daher
$Q_4=\tfrac16(R^2-3R_{\mu\nu}R^{\mu\nu})=\tfrac16(144-108)=6$. Mit Volumenelement
$\sqrt g=\sin^3\theta$, $\mathrm{Vol}(S^3)=2\pi^2$ und
\[
\int_0^{\pi/2}\sin^3\theta\,d\theta=\Bigl[-\cos\theta+\tfrac13\cos^3\theta\Bigr]_0^{\pi/2}=\tfrac23
\]
erhalten wir $\int_{D^4}Q_4\,\sqrt g\,d^4x=6\cdot 2\pi^2\cdot\tfrac23=8\pi^2.$
\end{proof}
\begin{remark}
$\CQ=8\pi^2$ ist exakt und maßstabsfrei. Das Verhältnis $\int_{D^4}Q_4/\int_{D^4}E_4=1/4$
(Ergebnis $R12$ unten) fixiert die relative $Q$-zu-$E$-Normierung.
\end{remark}

\section{Spektraler Unterbau: die zwölf Ergebnisse und die BFK-Kontrolle}\label{sec:spectral}
Die Konstruktion ruht auf zwölf $\zeta$-Funktions-Ergebnissen auf $S^4$, ihrer Hemisphäre $D^4$ und
dem äquatorialen $S^3$. Sie werden knapp angegeben; $R1$ ist das tragende und wird zweimal
hergeleitet.

\begin{description}[leftmargin=2.6em,nosep,itemsep=2pt]
\item[R1.] $\zeta(0;\mathrm{Maxwell};S^4)=-4\ag=-31/45$. Zwei Herleitungen stimmen überein: der
integrierte Spuranomalie-Weg ($-\ag\int E_4/16\pi^2=-4\ag$) und die Spektralkombination
\eqref{eq:zetacombo}. Jede Einheit der Euler-Charakteristik trägt $-2\ag$ bei; $\chi(S^4)=2$.
\item[R2.] $\zeta(0;\mathrm{DtN};\mathrm{Skalar};S^3)=-1$. Der skalare
Dirichlet-zu-Neumann-Operator auf dem äquatorialen $S^3$ (harmonische Fortsetzung in die Kappe) hat
Spektrum $\lambda_\ell=\ell(\ell+2)/(\ell+1)$, Multiplizität $(\ell+1)^2$; die analytische
Fortsetzung gibt $\zeta(0)=\zeta_R(-2)-1=-1$.
\item[R3.] $\zeta(0;\mathrm{DtN};\text{1-Form};S^3)=1$.
\item[R4.] $\zeta(0;\mathrm{DtN};\mathrm{Maxwell};S^3)=3=\dim S^3$ (eichfixierte Maxwell-DtN).
\item[R5.] $\zeta(0;\mathrm{Maxwell};S^3)=1$ (äquatoriale Photon-Identität).
\item[R6.] BFK-Identität: mit $R1$ und den beiden Buchhaltungen unten gilt
$2\zeta(0;D^4)-\zeta(0;\mathrm{DtN})=-31/45$.
\item[R7.] Photon-Casimir-Energie auf $S^3$: $E_{\mathrm{Cas}}=11/120$.
\item[R8.] $\det{}'\Lambda^{\mathrm{Maxwell}}_{S^3}=1/(2\pi^3)$.
\item[R9.] $\eta$ der transversalen $1$-Form-DtN ist gleich ihrem $\zeta(0)=1$, konform invariant
(verschieden von der de-Rham-APS-$\eta_{\mathrm{APS}}(S^3)=0$ der topologischen Quantisierung).
\item[R10.] Eine Robin-Bedingung $(\partial_n+\alpha)u=0$ am Äquator verschiebt das DtN-Spektrum um
$\alpha=\tan\ag\approx\ag$ --- der analytische Ursprung des Cap-Sattel-Antriebs (\S\ref{sec:cap}).
\item[R11.] Unter der Robin-Verschiebung ist $\zeta(0;\Lambda_{\mathrm{Robin}})$ linear in $\alpha$.
\item[R12.] $\int_{D^4}Q_4/\int_{D^4}E_4=1/4$ exakt.
\end{description}

\paragraph{Zwei konsistente BFK-Buchhaltungen.} Die Burghelea--Friedlander--Kappeler-Verklebungs\-identität
verbindet die $\zeta$-Determinante eines Operators auf einer geschlossenen Mannigfaltigkeit mit denen
auf zwei Teilen plus einem Dirichlet-zu-Neumann-Randterm. Zerlegt man $S^4$ in zwei Hemisphären, so
wird der geschlossene Wert $R1$ auf zwei intern konsistente Weisen wiedergewonnen, die sich nur in
der Eich-/Geist-Zuordnung unterscheiden:
\begin{align}
\text{(a) Skalar-DtN-Aufteilung:}\quad
&\zeta(0;D^4)=-2\ag-\tfrac12=-\tfrac{38}{45},
&&2\bigl(-\tfrac{38}{45}\bigr)-(-1)=-\tfrac{31}{45}\ \checkmark\\
\text{(b) eichfixierte Aufteilung:}\quad
&\zeta(0;D^4)=\tfrac{52}{45},
&&2\bigl(\tfrac{52}{45}\bigr)-3=-\tfrac{31}{45}\ \checkmark
\end{align}
Ihre Übereinstimmung ist eine interne Kontrolle der Spektralarithmetik; $R1$ selbst wird in
\S\ref{sec:heat} unabhängig aus dem Wärmekern \emph{hergeleitet}.

\section{Das Hierarchie-Gesicht: Sequestrierung auf der Hemisphäre}\label{sec:seq}
Eine naive Vakuumenergie skaliert wie $\Mpl^4$. Der Kaloper--Padilla-\emph{Sequestrierungs}-Mechanismus
macht die kosmologische Konstante zu einem globalen Lagrange-Multiplikator statt zu einer lokalen
Quelle: eine globale Zwangsbedingung entfernt den strahlungsinstabilen Bulk-Beitrag und lässt nur
das Rand-/topologische Stück übrig. Auf der Hemisphäre bleibt genau der Anomalie-Term: mit
$\int_{D^4}=\tfrac12\int_{S^4}$,
\begin{equation}\label{eq:sequester}
\tfrac12\,\Lambda_{\mathrm{eff}}\,V_4(D^4)=\ag .
\end{equation}
Entscheidend ist die Struktur: nach der Sequestrierung ist der verbleibende kosmologische Term durch
die \emph{dimensionslose} Anomalie-Ladung fixiert, nicht durch den Cutoff. Das macht ein
parameterfreies Verhältnis erst möglich und adressiert das Hierarchie-Gesicht
(\S\ref{sec:problem}, (i)).

\section{Der Infrarot-Lift durch kohärente spektrale Akkumulation}\label{sec:lift}
Die lokale Anomalie-Dichte auf dem de-Sitter-Raum ist
$\rho_{\mathrm{anom}}=\tfrac14\avg{T^\mu{}_\mu}=\tfrac{\ag}{8\pi^2}\,3H^4$, also
$\rho_{\mathrm{anom}}/\rho_{\mathrm{krit}}=\tfrac{\ag}{8\pi^2}(H/\Mpl)^2\sim10^{-122}$ --- eine rein
lokale Lesart ist viel zu klein. Das Schließen dieser Lücke ist eine Eigenschaft des Spektrums, keine
zusätzliche Annahme.

\begin{proposition}[Kohärente Akkumulation]\label{prop:lift}
Der eichfixierte Maxwell-Operator auf $S^4$ hat das diskrete Spektrum $\lambda_n=n(n+2)H^2$ mit
Multiplizität $d_n=2n(n+2)$, $n\ge1$ \cite{Allen,Higuchi}. Diese Moden sind durch die globale
Geometrie fixiert und phasengekoppelt, sodass die $\zeta$-regularisierte Vakuumsumme \emph{kohärent}
ist (linear in der Modenzahl $N$) und nicht stochastisch ($\sqrt N$). Die Bulk-Zählung bis zum Cutoff
$n_{\max}=\Mpl/H$ ist
\[
N_{\mathrm{tot}}=\sum_{n=1}^{n_{\max}}2n(n+2)=\tfrac23 n_{\max}^3+3n_{\max}^2+\tfrac73 n_{\max}
\sim\tfrac23\Bigl(\frac{\Mpl}{H}\Bigr)^3,
\]
und die Horizont-Flächen-Projektion reduziert dies auf
$N_{\mathrm{eff}}\sim(\Mpl/H)^2=S_{dS}/\pi$, die Gibbons--Hawking-de-Sitter-Entropie. Die
$(\Mpl/H)^2$-Skalierung ist exakt; ihr $O(1)$-Vorfaktor ist konventionsabhängig.
\end{proposition}

\begin{corollary}
$N_{\mathrm{eff}}\,\rho_{\mathrm{anom}}/\rho_{\mathrm{krit}}=\ag/8\pi^2$, also das dimensionslose
Verhältnis aus \S\ref{sec:friedmann}. Die Größe des \emph{Verhältnisses} --- der tatsächlich
gemessenen Größe --- ist damit ohne freien Parameter fixiert.
\end{corollary}

\section{Die integrierte Invariante $4\ag$ über zwei unabhängige Wege}\label{sec:invariant}
\begin{theorem}[Spektraler Weg]\label{thm:zeta}
$\displaystyle\int_{S^4}\avg{T^\mu{}_\mu}=4\ag=\bigl\lvert\zeta(0;\mathrm{Maxwell};S^4)\bigr\rvert.$
\end{theorem}
\begin{proof}
Die konforme Anomalie ist die Variation der Ein-Schleifen-Wirkung unter konstantem Weyl,
$\delta_\sigma\Gamma=-\tfrac12\zeta(0)\delta\sigma$, und $\zeta(0)$ ist das integrierte $a_4$. Nach
$R1$ gilt $\zeta(0;\mathrm{Maxwell};S^4)=-31/45=-4\ag$.
\end{proof}

\begin{theorem}[Skalenfluss-Weg]\label{thm:scaleflow}
Unter $g\to e^{2\sigma}g$ erfüllt die anomalie-induzierte (Riegert--Mottola-)Wirkung die
Wess--Zumino-Konsistenzrelation mit
$\displaystyle\frac{dS_{\mathrm{anom}}}{d\sigma}=\frac{\ag}{16\pi^2}\int_{S^4}E_4=4\ag,$
erneut wegen $\int_{S^4}E_4=64\pi^2$. \cite{Riegert1984,Mottola}
\end{theorem}

\begin{remark}
Die Sätze~\ref{thm:zeta} und~\ref{thm:scaleflow} fixieren die Zahl $4\ag$ über zwei Konstruktionen
ohne gemeinsamen Zwischenschritt und ohne freien $O(1)$-Faktor. Der Wert $4\ag$ ist daher
vorfaktorunabhängig.
\end{remark}

\section{Der Variationsmechanismus, Stabilität und das Vorzeichen}\label{sec:cap}
\begin{proposition}[Cap-Sattel und das offene Universum]\label{prop:sign}
Ergebnis $R10$ zeigt, dass eine Robin-Bedingung am Äquator das DtN-Spektrum um $\alpha=\tan\ag$
verschiebt. Die Extremierung der resultierenden Rand-Wirkung über die äquatoriale Lage
$\delta\theta$ ergibt die Stationaritätsgleichung
\[
\ag\cos^4(\delta\theta)=\sin(\delta\theta),
\]
mit einer stabilen Lösung $\delta\theta^*\approx0.164>0$. Die Stabilität folgt aus dem positiven
diskreten Spektrum des Laplace-Operators auf der kompakten $S^4$ (ein Maximum der euklidischen
Wirkung ist nach der Wick-Rotation ein Minimum der Energie). Das positive Vorzeichen von
$\delta\theta^*$ erzwingt $\Ok>0$: ein \emph{offenes} Universum. Die Ladung $\CQ=8\pi^2$ wird stets
auf der Standard-Hemisphäre ausgewertet, sodass das führende Verhältnis
$\lvert\Ok\rvert/\OL=\ag/8\pi^2$ keine $\delta\theta^*$-Korrektur trägt; der Cap-Sattel geht erst auf
der Verfeinerungsebene ein.
\end{proposition}

\section{Eindeutigkeit: warum das Photon}\label{sec:uniq}
Wir halten fest, warum kein anderes Feld das Photon ersetzt und warum kein alternativer Mechanismus
den Absolutwert schließt. \emph{(a) Feldinhalt:} massive Felder entkoppeln wie $m^2/\Mpl^2$; unter
den masselosen Feldern geht der Typ-A-Koeffizient des Gravitons nur mit propagierender Quantengravitation
ein, die wir nicht heranziehen, während Neutrinos (falls masselos) und Gluonen entweder nicht im
relevanten Sinn konform gekoppelt oder eingeschlossen sind; das ungebrochene $U(1)$-Photon ist der
saubere Beitrag. \emph{(b) Mechanismus:} kandidierende dynamische Schließungen des Absolutwerts (ein
Scalaron mit Coleman--Weinberg-Potential; die Identifikation eines zweiten Operators mit Dunkler
Materie; eine stochastische statt kohärente Summe) wurden jeweils geprüft und liefern entweder die
falsche Zahl ($1/(12\ag)\approx0.48$ statt $0.69$), sind lokal vernachlässigbar oder verletzen die in
Proposition~\ref{prop:lift} gezeigte Kohärenz. Die überlebende Aussage ist das bewiesene
\emph{Verhältnis} zusammen mit dem einzigen identifizierenden Postulat aus \S\ref{sec:postulate}.

\section{Das parameterfreie Verhältnis und die Friedmann-Übersetzung}\label{sec:friedmann}
\begin{theorem}[Zentrale Beziehung]\label{thm:central}
Mit $\CQ=8\pi^2$ (Lemma~\ref{lem:qcharge}) und $\ag=31/180$ (Satz~\ref{thm:agamma}) gilt
\begin{equation}\label{eq:central}
\frac{\lvert\Ok\rvert}{\OL}=\frac{\ag}{\CQ}=\frac{\ag}{8\pi^2}=\frac{31}{1440\pi^2}\approx2.181\times10^{-3}.
\end{equation}
Die Beziehung enthält keinen Maßstab, keine Newton-Konstante, keine Planck-Masse.
\end{theorem}

\paragraph{Acht-Schritte-Zusammenbau.}
\emph{(1)} $\int_{S^4}\avg{T}=4\ag$ (Satz~\ref{thm:gbc}); \emph{(2)} Sequestrierung auf $D^4$,
Gl.~\eqref{eq:sequester}; \emph{(3)} $\CQ=8\pi^2$ (Lemma~\ref{lem:qcharge}); \emph{(4)}
Wess--Zumino-Linearantwort (Satz~\ref{thm:scaleflow}); \emph{(5)} der dimensionslose Quotient
\eqref{eq:central}; \emph{(6)} Photon-Auswahl (\S\ref{sec:uniq}); \emph{(7)} der euklidische
Cap-Sattel (Proposition~\ref{prop:sign}) und seine Wick-Rotation in einen Lorentzschen
Friedmann-Patch; \emph{(8)} die Friedmann-Übersetzung unten. Schritte $1$--$6$ und die euklidischen
Teilschritte von $7$ sind satz-niveau; das einzige genuin gewählte Element ist die
Gegenwarts-Identifikation, die in \S\ref{sec:postulate} isoliert wird.

\paragraph{Friedmann-Übersetzung.} Mit $\overline\Lambda=8\pi^2/\ag=1440\pi^2/31\approx458.4$ fixiert
das Postulat $\OL=4\ag$ (\S\ref{sec:postulate}) zusammen mit $\OL=\overline\Lambda\lvert\Ok\rvert$
\begin{equation}\label{eq:OK}
\lvert\Ok\rvert=\frac{\OL}{\overline\Lambda}=\frac{\ag^2}{2\pi^2}\approx1.50\times10^{-3}
\quad(\text{offen, }\Ok>0),
\end{equation}
und das Flachheitsbudget $\Om+\OL+\Ok=1$ \emph{sagt} dann $\Om=1-\OL-\lvert\Ok\rvert\approx0.3096$
\emph{voraus}, verträglich mit dem Planck-2018-Wert $\Om=0.3111\pm0.0056$; äquivalent stimmt
$\OL=0.6889$ mit Planck auf vier Dezimalen überein \cite{Planck2018}.

\section{Das einzige physikalische Postulat}\label{sec:postulate}
Alles in \S\S\ref{sec:anomaly}--\ref{sec:friedmann} ist euklidisch und geometrisch. Ein Schritt
verbindet den bewiesenen euklidischen Quotienten mit dem heutigen Friedmann-Observablen.

\begin{postulate}[Anomalie--Krümmung-Korrespondenz, ``A5/P5'']\label{post:A5}
Das auf der euklidischen Hemisphäre berechnete Anomalie-Verhältnis $\ag/\CQ$, durch die
Wick-Rotation in einen Lorentzschen Friedmann-Patch getragen, ist gleich dem beobachteten heutigen
Parameter der räumlichen Krümmung; äquivalent sättigt die Dunkle-Energie-Dichte die integrierte
Invariante, $\OL=4\ag$.
\end{postulate}

\begin{remark}[Genauer Status von Postulat~\ref{post:A5}]
Die euklidische Hälfte --- die topologische Identität $\Ok^{\mathrm{konform}}=\ag/\CQ$ --- ist ein
\emph{Satz}: die Cauchysche Funktionalgleichung liefert die Linearität, ein Chern--Weil-Basiswechsel
das $1/8\pi^2$, und der Atiyah--Patodi--Singer-Indexsatz auf dem de-Rham-Komplex die ganzzahlige
Normierung (bis auf eine beschränkte Randverfeinerung
$\lvert\epsilon_\partial\rvert\le1.5\%$). Die Wick-Rotations-Kinematik ist in der
Sattelpunktnäherung \emph{herleitbar}. Was eine genuine \emph{physikalische Hypothese} bleibt, ist
die Identifikation der resultierenden euklidischen Größe mit dem heutigen Observablen unter der
kanonischen $\Lambda$-dominierten Referenzepoche. Sie ist physikalisch motiviert --- $4\ag$ ist
zugleich die spektrale $\zeta$-Invariante (Satz~\ref{thm:zeta}) und der
Weyl-Skalenanomalie-Koeffizient (Satz~\ref{thm:scaleflow}) ---, aber sie wird \emph{nicht}
dynamisch erzwungen: die Anomalie-Wirkung $S_{\mathrm{anom}}(\sigma)=4\ag\,\sigma$ ist linear in
$\sigma$ und hat keinen stationären Punkt, und der dynamische Scalaron-Weg liefert
$1/(12\ag)\approx0.48$ statt $0.69$. Wir stellen sie daher als Gründungs-Postulat auf, von derselben
logischen Art wie die postulierten Prinzipien der etablierten physikalischen Theorien
(Äquivalenzprinzip, Schrödinger-Gleichung, Born-Regel, Eichprinzip), mit dem unterscheidenden
Merkmal, dass ihr Wert durch die zwei Invarianten aus \S\ref{sec:invariant} unabhängig fixiert
ist. Sie wird nicht als hergeleiteter Satz behauptet und ist der eigentliche Angriffspunkt jedes
Falsifizierungsversuchs.
\end{remark}

\section{Vorhersagen, Beobachtungsstand, Prognosen}\label{sec:obs}
Gegeben Postulat~\ref{post:A5}:
\[
\OL=4\ag\approx0.6889,\qquad \Ok\simeq+1.5\times10^{-3}\ (\text{offen}),\qquad w=-1 .
\]
Das Verhältnis $\lvert\Ok\rvert/\OL=\ag/8\pi^2$ aus Satz~\ref{thm:central} gilt \emph{unabhängig}
vom Postulat. Alle sind falsifizierbar: ein gemessenes $w$, das bei hoher Signifikanz von $-1$
abweicht, ein geschlossenes Universum ($\Ok<0$) oder ein Krümmungs-zu-Dunkle-Energie-Verhältnis, das
von $2.181\times10^{-3}$ abweicht, widerlegt jeweils das Rahmenwerk. Der
Zustandsgleichungswert $w=-1$ steht derzeit in $\sim3$--$4\sigma$-Spannung zu DESI~DR2; DESI~DR3 und
Euclid werden entscheiden. Das ist eine scharfe, riskante Vorhersage, keine Rückschau.

\section{Logischer Status jeder Aussage}\label{sec:ledger}
\begin{center}
\begin{tabular}{lll}
\toprule
Aussage & Status & Methode / Referenz \\
\midrule
$\ag=31/180$ & bewiesen & Gilkey $a_4$, drei Operatoren (Satz~\ref{thm:agamma}) \\
$\int_{S^4}E_4=64\pi^2$; Vorfaktor $4=2\chi$ & bewiesen & Gauß--Bonnet--Chern (Satz~\ref{thm:gbc}) \\
$\CQ=8\pi^2$ & bewiesen & elementares Integral (Lemma~\ref{lem:qcharge}) \\
spektraler Unterbau $R1$--$R12$ & bewiesen & \S\ref{sec:spectral}; $R1$ zweifach \\
$\int_{S^4}\avg T=4\ag$ & bewiesen (zwei Wege) & $\zeta$-Determinante; Skalenfluss (Sätze~\ref{thm:zeta},~\ref{thm:scaleflow}) \\
Verhältnis $\ag/8\pi^2$ & bewiesen, parameterfrei & Satz~\ref{thm:central} \\
Sequestrierungs-Unterdrückung & Literatur-Mechanismus & Kaloper--Padilla (\S\ref{sec:seq}) \\
Lift-$(\Mpl/H)^2$-Skalierung & bewiesen (Skalierung) & kohärente Spektralsumme (Prop.~\ref{prop:lift}) \\
Vorzeichen $\Ok>0$ & bewiesen & spektrale Positivität (Prop.~\ref{prop:sign}) \\
Photon-Eindeutigkeit & argumentiert (Ausschöpfung) & \S\ref{sec:uniq} \\
$\OL=4\ag$ (Absolutwert) & \textbf{postuliert} & Identifikation A5/P5 (Post.~\ref{post:A5}) \\
numerische $\Om,\OL,\Ok$; $w=-1$ & empirisch & DESI~DR3 / Euclid (offen) \\
\bottomrule
\end{tabular}
\end{center}

\noindent
Jede mit ``bewiesen'' markierte Zeile wird zur unabhängigen Überprüfung für sich angeboten. Die
einzige tragende Hypothese ist die mit \textbf{postuliert} markierte Zeile; jede Falsifizierung des
Rahmenwerks sollte dort oder an den empirischen Vorhersagen aus \S\ref{sec:obs} ansetzen.

\begin{thebibliography}{99}
\bibitem{Gilkey} P.~B.~Gilkey, \emph{Invariance Theory, the Heat Equation, and the
Atiyah--Singer Index Theorem}, 2.~Aufl., CRC Press (1995).
\bibitem{Birrell1982} N.~D.~Birrell, P.~C.~W.~Davies, \emph{Quantum Fields in Curved Space},
Cambridge (1982).
\bibitem{ChristensenDuff1980} S.~M.~Christensen, M.~J.~Duff, Nucl. Phys. B \textbf{170} (1980) 480.
\bibitem{Vassilevich2003} D.~V.~Vassilevich, \emph{Heat kernel expansion: user's manual}, Phys.
Rep. \textbf{388} (2003) 279.
\bibitem{Duff1994} M.~J.~Duff, \emph{Twenty years of the Weyl anomaly}, Class. Quantum Grav.
\textbf{11} (1994) 1387.
\bibitem{DeserSchwimmer} S.~Deser, A.~Schwimmer, Phys. Lett. B \textbf{309} (1993) 279.
\bibitem{Chern} S.-S.~Chern, Ann. Math. \textbf{45} (1944) 747.
\bibitem{AtiyahSinger} M.~F.~Atiyah, I.~M.~Singer, Ann. Math. \textbf{87} (1968) 484.
\bibitem{Riegert1984} R.~J.~Riegert, Phys. Lett. B \textbf{134} (1984) 56.
\bibitem{Mottola} E.~Mottola, anomalie-induzierte effektive Wirkung (und Referenzen darin).
\bibitem{Allen} B.~Allen, Vektor-Zweipunktfunktionen auf maximal symmetrischen Räumen (1986).
\bibitem{Higuchi} A.~Higuchi, J. Math. Phys. \textbf{28} (1987) 1553.
\bibitem{BGT1995} M.~Bucher, A.~Goldhaber, N.~Turok, Phys. Rev. D \textbf{52} (1995) 3314.
\bibitem{HawkingTurok1998} S.~W.~Hawking, N.~Turok, Phys. Lett. B \textbf{425} (1998) 25.
\bibitem{Planck2018} Planck Collaboration, \emph{Planck 2018 results. VI}, A\&A \textbf{641} (2020) A6.
\end{thebibliography}

\end{document}
